\input{configTD}
\usepackage{listings}

\lstset{
    basicstyle=\ttfamily\small,
    breaklines=true,
    frame=single,
    columns=fullflexible,
    keepspaces=true
}

\title{Analyse Numérique TP2 - Systèmes linéaires, conditionnement et comparaison de méthodes}

\author{}
\date{}

\begin{document}
\sffamily
\maketitle

\section*{Objectifs}
\begin{itemize}
    \item Comparer une méthode directe et une méthode itérative pour résoudre un système linéaire.
    \item Mesurer le temps de calcul sur des matrices de tailles différentes.
    \item Observer l'effet d'une perturbation du second membre.
    \item Relier une matrice tridiagonale à une discrétisation de type chaleur ou barre élastique.
\end{itemize}

\section*{Partie 1 - Une matrice issue d'une discrétisation}
On considere pour $n\geq 2$ la matrice tridiagonale
\[
A_n=
\begin{pmatrix}
2 & -1 & 0 & \cdots & 0\\
-1 & 2 & -1 & \ddots & \vdots\\
0 & -1 & 2 & \ddots & 0\\
\vdots & \ddots & \ddots & \ddots & -1\\
0 & \cdots & 0 & -1 & 2
\end{pmatrix}.
\]
Cette matrice apparaît dans plusieurs problèmes de différences finies ou d'éléments finis 1D.

\begin{enumerate}
    \item Écrire une fonction Python qui construit $A_n$.

    \item Construire le second membre
    \[
    b_n=(1,\dots,1)^T.
    \]

    \item Vérifier pour quelques tailles que $A_n$ est inversible.
\end{enumerate}

\begin{lstlisting}[language=Python]
import numpy as np
import time

def build_matrix(n):
    A = 2 * np.eye(n)
    A += -1 * np.eye(n, k=1)
    A += -1 * np.eye(n, k=-1)
    return A

def build_rhs(n):
    return np.ones(n)
\end{lstlisting}

\section*{Partie 2 - Méthode directe}
On utilisera d'abord la résolution directe de \texttt{numpy}.

\begin{enumerate}
    \item Résoudre $A_nx=b_n$ pour $n=10$, $100$ et $500$.

    \item Mesurer le temps de calcul avec \texttt{time.perf\_counter()}.

    \item Tracer le temps de calcul en fonction de $n$.
\end{enumerate}

\begin{lstlisting}[language=Python]
def solve_direct(A, b):
    return np.linalg.solve(A, b)

sizes = [10, 100, 500]
times_direct = []

for n in sizes:
    A = build_matrix(n)
    b = build_rhs(n)
    t0 = time.perf_counter()
    x = solve_direct(A, b)
    t1 = time.perf_counter()
    times_direct.append(t1 - t0)
    print(n, times_direct[-1], np.linalg.norm(A @ x - b))
\end{lstlisting}

\section*{Partie 3 - Méthode de Jacobi}
On implémente maintenant la méthode itérative de Jacobi.

\begin{enumerate}
    \item Compléter la fonction suivante.

    \item Pour $n=10$, faire afficher le nombre d'itérations nécessaires pour atteindre
    \[
    \|x^{(k+1)}-x^{(k)}\|_{\infty} < 10^{-8}.
    \]

    \item Recommencer pour $n=100$ puis $n=500$.
\end{enumerate}

\begin{lstlisting}[language=Python]
def jacobi(A, b, x0=None, tol=1e-8, nmax=20000):
    n = len(b)
    if x0 is None:
        x = np.zeros(n)
    else:
        x = x0.copy()

    D = np.diag(A)
    R = A - np.diag(D)

    for k in range(nmax):
        x_new = (b - R @ x) / D
        if np.linalg.norm(x_new - x, ord=np.inf) < tol:
            return x_new, k + 1
        x = x_new

    return x, nmax
\end{lstlisting}

\section*{Partie 4 - Comparaison directe / itérative}
\begin{enumerate}
    \item Comparer les temps de calcul des deux approches.

    \item Comparer les résidus
    \[
    \|Ax-b\|_{\infty}
    \]
    obtenus par les deux méthodes.

    \item Discuter l'intérêt pratique d'une méthode itérative lorsque la matrice est grande et creuse.
\end{enumerate}

\section*{Partie 5 - Perturbation du second membre}
On modifie maintenant le second membre en posant
\[
\tilde b_n=b_n+\varepsilon e_1
\]
ou $e_1=(1,0,\dots,0)^T$ et $\varepsilon=10^{-3}$.

\begin{enumerate}
    \item Résoudre $A_n\tilde x=\tilde b_n$.

    \item Comparer $x$ et $\tilde x$ avec la norme infinie.

    \item Calculer le rapport
    \[
    \frac{\|\tilde x-x\|_{\infty}}{\|x\|_{\infty}}
    \qquad \text{et} \qquad
    \frac{\|\tilde b-b\|_{\infty}}{\|b\|_{\infty}}.
    \]

    \item Interpréter le résultat.
\end{enumerate}

\begin{lstlisting}[language=Python]
def perturb_rhs(b, eps=1e-3):
    bp = b.copy()
    bp[0] += eps
    return bp
\end{lstlisting}

\section*{Partie 6 - Nombre de condition}
\begin{enumerate}
    \item Calculer \texttt{np.linalg.cond(A)} pour plusieurs tailles.

    \item Commenter l'évolution du conditionnement avec la taille.

    \item Relier cette observation à la sensibilité des calculs précédents.
\end{enumerate}

\section*{Partie 7 - Bonus applicatif}
Choisir l'une des deux ouvertures suivantes.

\subsection*{Option A - PageRank}
Construire une petite matrice de transition stochastique associée à un graphe de $4$ pages, puis approcher sa distribution stationnaire par itérations.

\subsection*{Option B - Element fini 1D}
Interpréter $A_n$ comme matrice de rigidité d'une barre 1D et commenter le lien entre maillage fin, taille de la matrice et coût numérique.

\section*{Compte rendu attendu}
\begin{itemize}
    \item Les fonctions Python sont exécutées.
    \item Les temps de calcul sont comparés sur plusieurs tailles.
    \item L'effet de la perturbation est discuté.
    \item Le rôle du conditionnement est expliqué en quelques lignes.
    \item Le bonus est traité au moins partiellement.
\end{itemize}

\end{document}
